\section{Selective Recovery Model}
\label{sec:recovery-model}

Rollback wastes valid work when its boundary follows execution order rather
than dependency. As Section~\ref{sec:motivation} shows, affected results and
independent work can interleave within a turn, checkpoint interval, or branch.
We therefore separate the speculative execution unit from the recovery unit.
The complete trajectory shares speculative state; a dependency subgraph
determines which effects recovery removes.

\subsection{Overview and Goals}
\label{sec:aet-model}

We describe this separation as an \emph{Agent Effect Transaction} (AET).
An AET records tool steps, their observed effects, historical object states,
and associated agent-visible messages throughout execution.
Figure~\ref{fig:overview} shows the workflow: capture evidence, select a
candidate recovery set, reconstruct affected state, and continue with
independent work retained.
Section~\ref{sec:design} describes the Linux implementation.

\begin{figure}[t]
\centering
\includegraphics[width=\linewidth]{FIG-AgentTX-Overview.pdf}
\caption{\textbf{Execution and selective recovery in \system.}
\small\emph{Recorded dependencies determine the invalidation set.
If a recovery plan selects E3, its descendants E4--E7 are also invalidated,
while independent E8 remains in the workspace and visible context.}}
\label{fig:overview}
\end{figure}

Consider an interleaved trajectory: $s_0$ modifies a formatter, $s_1$ uses it
to generate a report, $s_2$ edits independent documentation, and $s_3$ checks
the report.
If a recovery plan selects $s_0$, retaining the report would retain its
affected result, whereas restoring an earlier checkpoint would discard the
documentation.
The dependency graph can select $s_0$, $s_1$, and $s_3$ while preserving $s_2$.

The design has three goals corresponding to the preceding challenges.
First, discover recovery dependencies from observed execution and result
flow, so invalid results are removed without automatically discarding later
work.
Second, associate effects with objects and versions consistently across
dependency discovery and restoration.
Third, reconstruct a non-prefix workspace and visible context that describe
the same surviving steps.

\subsection{An Observed Dependency Graph}
\label{sec:effect-dag}

The recovery graph represents tool steps as nodes and observed recovery
dependencies as forward edges.
An edge connects a producer to a later consumer when the latter reads
state produced by the former.
The graph also records conservative constraints from overlapping mutations,
deletions, and namespace observations, including failed lookups followed by
creation.
These constraints prevent recovery from separating effects whose recorded
state relationships require coordinated removal.

Object association determines which accesses belong to these relationships.
Different names may identify the same object, while replacement can make one
name refer to different objects over time.
The AET therefore relates observed accesses to supported object identities
and historical states, rather than relying solely on path equality.
These associations serve both graph construction and later version selection.

Dependencies can also pass through tool results without an intervening file.
When a later request reuses a detectable value from an earlier result, the
graph connects their steps.
Filesystem and message edges form a graph that can branch and merge; the
causal chain of invalidation need not follow one chronological sequence.

\subsection{From Candidates to an Invalidation Set}
\label{sec:recovery-selection}

Failure observation and recovery-set selection answer different questions.
An observed failure motivates candidate changes to withdraw; the graph
determines which descendants must accompany each candidate.
Candidate localization and isolated verification are described separately
in Section~\ref{sec:design}.

For a candidate set $B$, recovery begins with $C=B$ and repeatedly adds
dependent steps until the set stops growing.
Expansion follows both filesystem and message edges, including paths that
cross between the two channels.
Invalidation need not be contiguous: membership depends on reachability,
so independent steps between selected nodes remain eligible to survive.
The surviving state therefore need not be a temporal prefix.
Figure~\ref{fig:recovery-workflow} summarizes this selection and the
reconstruction guard.

\begin{figure}[t]
\centering
\begin{lstlisting}[style=agenttxcode]
// Input: candidate steps B; effects and history.
C = B;
do {
    old = C;
    C |= file_descendants(C, effects);
    C |= message_descendants(C, history);
} while (C != old);

objects = mutated_objects(C, effects);
// Versions before the earliest selected mutations.
versions = prior_versions(objects, C, history);
// Check identity, history, and retained mutations.
if (!can_restore(objects, versions, C, effects))
    return unsupported_plan;
return plan(C, versions); // For isolated checking.
\end{lstlisting}
\caption{\textbf{Candidate-derived recovery workflow.}
\small\emph{Joint dependency closure determines the invalidation set;
reconstruction checks determine whether it can be applied safely.}}
\label{fig:recovery-workflow}
\end{figure}

The fixed point is the least set containing the candidates and closed under
the recorded filesystem and message dependencies.

\subsection{Reconstructing Affected Object Versions}
\label{sec:selective-reconstruction}

Selecting invalid steps does not yet produce a usable workspace.
Repeated writes and deletions can hide earlier valid object versions, so
deleting the latest files alone cannot undo selected effects.
Recovery instead selects historical state preceding the earliest selected
mutation of each affected object group.
It checks these targets against retained mutations before applying the
plan, rejecting unsupported overlaps rather than guessing a merge.

Historical restoration changes affected objects while independent state
remains in place.
The result may combine versions that never coexisted at any chronological
checkpoint.
After successful reconstruction, selected steps become inactive in the
execution history.
Their mutations are absent and nonconflicting retained effects remain visible.

\subsection{Continuing with Matching Visible Context}
\label{sec:context-reconstruction}

The same invalidation set governs workspace and visible-context
reconstruction.
Otherwise, an agent could continue using a tool result whose underlying
state recovery has already withdrawn.
Recovery removes selected tool-call spans and stale recovery-control
messages, retains independent spans, and describes the surviving work.
In the example, the documentation and its valid record survive together;
the affected report and its associated results are invalidated together.

The next model request therefore describes the same survivor set as the
reconstructed workspace.
Native integrations rebuild the request from retained spans; black-box
agents receive a retained-effects manifest instead.
The agent then continues repair and verification from the retained
progress.
